\documentclass[../DoAn.tex]{subfiles}
\begin{document}

Chương trước đã trình bày phương pháp đề xuất của đồ án theo hướng hậu xử lý kết quả ước lượng tư thế người 3D từ hai video monocular. Hướng tiếp cận này không thay thế mô hình ước lượng tư thế ban đầu, mà sử dụng kết quả từ WHAM làm đầu vào để tiếp tục đồng bộ, căn chỉnh, kiểm tra nhất quán, phân xử và hiệu chỉnh pose \cite{shin2024wham}. WHAM là phương pháp tái tạo chuyển động người 3D từ video monocular trong hệ tọa độ toàn cục, do đó phù hợp để tạo chuỗi pose thô ban đầu cho pipeline hậu xử lý của đồ án \cite{shin2024wham}.

Chương này trình bày phân tích lý thuyết cho phương pháp đề xuất. Nội dung tập trung vào mô hình hóa bài toán refinement, phân tích các ràng buộc dữ liệu, thời gian, đa góc nhìn, động học/cơ sinh học, vai trò của SMPL và độ phức tạp tính toán của pipeline \cite{smpl2015,smplify2016,physcap2020}.

\section{Mô hình hóa bài toán refinement hậu xử lý}

Bài toán của đồ án được mô hình hóa là bài toán hậu xử lý trên chuỗi pose 3D đã được ước lượng sẵn, thay vì bài toán học trực tiếp ánh xạ từ ảnh RGB sang pose 3D \cite{shin2024wham}. Cách mô hình hóa này phù hợp vì các phương pháp monocular có thể tạo ra pose 3D từ một video, nhưng vẫn chịu ảnh hưởng của nhập nhằng chiều sâu, che khuất và thiếu ổn định theo thời gian \cite{smplify2016,physcap2020}. Đầu vào của hệ thống gồm hai chuỗi pose 3D được sinh độc lập từ hai video monocular của cùng một chuyển động \cite{shin2024wham}.

Gọi $X^{(1)}_{1:T}$ và $X^{(2)}_{1:T}$ lần lượt là hai chuỗi pose 3D tương ứng với camera thứ nhất và camera thứ hai. Với biểu diễn keypoints 3D, pose tại frame $t$ được mô tả bởi tập các khớp:

\begin{equation}
X_t = \{X_{t,j}\}_{j=1}^{J},
\end{equation}

trong đó $X_{t,j} \in \mathbb{R}^{3}$ là tọa độ 3D của khớp thứ $j$ tại frame $t$, $J$ là số lượng khớp và $T$ là số lượng frame. Biểu diễn pose dưới dạng tập khớp 3D là cách biểu diễn phổ biến trong bài toán 3D human pose estimation và cũng thuận tiện cho việc tính sai khác hình học giữa hai kết quả pose \cite{probtriangulation2023,gu2024multiview}.

Mục tiêu của refinement là tìm chuỗi pose đã hiệu chỉnh $\hat{X}_{1:T}$ sao cho kết quả này vừa gần với pose thô ban đầu, vừa thỏa mãn tốt hơn các ràng buộc về thời gian, đa góc nhìn và cấu trúc cơ thể \cite{smplify2016,physcap2020}. Bài toán có thể được viết dưới dạng tối ưu tổng quát như sau:

\begin{equation}
\hat{X}_{1:T}
=
\arg\min_{X_{1:T}}
\mathcal{L}_{total}.
\end{equation}

Trong đó, $\mathcal{L}_{total}$ là hàm mục tiêu tổng hợp nhiều thành phần sai số. Một dạng tổng quát của hàm mục tiêu được sử dụng trong phân tích của đồ án là:

\begin{equation}
\mathcal{L}_{total}
=
\lambda_{data}\mathcal{L}_{data}
+
\lambda_{temp}\mathcal{L}_{temp}
+
\lambda_{view}\mathcal{L}_{view}
+
\lambda_{bio}\mathcal{L}_{bio}
+
\lambda_{smpl}\mathcal{L}_{smpl}.
\end{equation}

Thành phần $\mathcal{L}_{data}$ giữ cho pose sau hiệu chỉnh không lệch quá xa so với output ban đầu của mô hình monocular \cite{smplify2016}. Thành phần $\mathcal{L}_{temp}$ biểu diễn ràng buộc nhất quán theo thời gian nhằm giảm jitter trong chuỗi chuyển động \cite{physcap2020,shin2024wham}. Thành phần $\mathcal{L}_{view}$ biểu diễn ràng buộc nhất quán giữa hai góc nhìn sau khi hai pose đã được đưa về cùng hệ quy chiếu \cite{gu2024multiview,probtriangulation2023}. Thành phần $\mathcal{L}_{bio}$ biểu diễn các ràng buộc động học và cơ sinh học nhằm hạn chế pose không hợp lý về cấu trúc cơ thể \cite{physcap2020,delp2007opensim,uhlrich2023opencap}. Thành phần $\mathcal{L}_{smpl}$ giữ nghiệm hiệu chỉnh trong không gian biểu diễn hợp lệ của mô hình SMPL \cite{smpl2015,smplify2016}.

\section{Phân tích ràng buộc dữ liệu}

Ràng buộc dữ liệu được sử dụng để bảo toàn thông tin từ pose thô ban đầu. Thành phần này cần thiết vì nếu chỉ tối ưu theo các ràng buộc làm mượt hoặc ràng buộc hình học, pose sau hiệu chỉnh có thể lệch khỏi chuyển động đã được mô hình monocular dự đoán từ video \cite{smplify2016,shin2024wham}. Với biểu diễn keypoints 3D, ràng buộc dữ liệu có thể được định nghĩa bằng sai số bình phương giữa pose sau refinement và pose khởi tạo:

\begin{equation}
\mathcal{L}_{data}
=
\sum_{t=1}^{T}
\sum_{j=1}^{J}
\left\|
X_{t,j} - X^{0}_{t,j}
\right\|_2^2.
\end{equation}

Trong đó, $X_{t,j}$ là vị trí khớp thứ $j$ tại frame $t$ sau refinement, còn $X^{0}_{t,j}$ là vị trí khớp tương ứng từ output ban đầu. Công thức này đóng vai trò như một regularization term, tương tự tinh thần của các phương pháp fitting khi cần cân bằng giữa dữ liệu quan sát và prior của mô hình cơ thể \cite{smplify2016,smpl2015}.

Nếu pose được biểu diễn trong không gian SMPL, ràng buộc dữ liệu có thể áp dụng lên tham số pose, shape và translation. SMPL là mô hình cơ thể người tham số, trong đó hình dạng và tư thế cơ thể được biểu diễn bằng các tham số có cấu trúc để sinh mesh cơ thể người \cite{smpl2015}. Khi đó, một ràng buộc trong không gian tham số có thể được viết như sau:

\begin{equation}
\mathcal{L}_{smpl}
=
\left\|
\theta - \theta^{0}
\right\|_2^2
+
\left\|
\beta - \beta^{0}
\right\|_2^2
+
\left\|
\tau - \tau^{0}
\right\|_2^2.
\end{equation}

Trong đó, $\theta$ là tham số pose, $\beta$ là tham số shape và $\tau$ là vector translation. Các giá trị có ký hiệu $0$ là kết quả khởi tạo từ mô hình đầu vào. Ràng buộc này phù hợp với hướng SMPLify, trong đó bài toán khôi phục cơ thể người được giải bằng cách tối ưu tham số SMPL sao cho phù hợp với tín hiệu quan sát và các prior của mô hình \cite{smplify2016,smpl2015}.

\section{Phân tích ràng buộc nhất quán theo thời gian}

Trong video, chuyển động cơ thể người thường có tính liên tục theo thời gian. Vì vậy, sự thay đổi quá đột ngột của các khớp giữa hai frame liên tiếp có thể là dấu hiệu của nhiễu dự đoán hoặc lỗi ước lượng pose \cite{physcap2020,shin2024wham}. Ràng buộc nhất quán theo thời gian được sử dụng để giảm jitter và làm ổn định quỹ đạo chuyển động của các khớp \cite{physcap2020}.

Một ràng buộc temporal đơn giản là sai phân bậc nhất giữa hai frame liên tiếp:

\begin{equation}
\mathcal{L}_{temp}^{(1)}
=
\sum_{t=2}^{T}
\sum_{j=1}^{J}
\left\|
X_{t,j} - X_{t-1,j}
\right\|_2^2.
\end{equation}

Thành phần này làm giảm biến đổi vị trí khớp giữa hai frame liền kề. Tuy nhiên, nếu trọng số của sai phân bậc nhất quá lớn, chuyển động có thể bị làm mượt quá mức và làm mất các động tác nhanh hợp lệ \cite{physcap2020}. Do đó, sai phân bậc hai thường phù hợp hơn khi mục tiêu chính là hạn chế biến thiên gia tốc bất thường:

\begin{equation}
\mathcal{L}_{temp}^{(2)}
=
\sum_{t=2}^{T-1}
\sum_{j=1}^{J}
\left\|
X_{t+1,j}
-
2X_{t,j}
+
X_{t-1,j}
\right\|_2^2.
\end{equation}

Sai phân bậc hai không ép hai vị trí khớp ở hai frame liên tiếp phải gần nhau tuyệt đối, mà chủ yếu hạn chế sự thay đổi đột ngột của vận tốc. Cách ràng buộc này phù hợp với mục tiêu giảm jitter trong chuỗi pose 3D, vì các phương pháp monocular có thể tạo ra kết quả thiếu ổn định theo thời gian khi gặp che khuất hoặc nhập nhằng chiều sâu \cite{shin2024wham,physcap2020}.

\section{Phân tích ràng buộc nhất quán đa góc nhìn}

Trong bài toán sử dụng hai camera monocular, cùng một chuyển động được quan sát từ hai góc nhìn khác nhau. Mỗi góc nhìn có thể sinh lỗi khác nhau do che khuất, góc nhìn bất lợi hoặc sai lệch chiều sâu trong ước lượng monocular \cite{smplify2016,physcap2020}. Vì vậy, việc so sánh hai kết quả pose độc lập có thể cung cấp tín hiệu để phát hiện các frame hoặc khớp có khả năng sai lệch \cite{gu2024multiview,probtriangulation2023}.

Trước khi so sánh hai pose, cần đưa chúng về cùng một hệ tọa độ hoặc một biểu diễn tương thích. Một cách tiếp cận phổ biến là sử dụng similarity transformation hoặc Procrustes alignment để căn chỉnh các tập điểm 3D \cite{gu2024multiview}. Phép căn chỉnh từ view thứ hai sang không gian so sánh chung có thể được viết như sau:

\begin{equation}
\tilde{X}^{(2)}_{t,j}
=
s_t R_t X^{(2)}_{t,j} + d_t.
\end{equation}

Trong đó, $s_t$ là hệ số tỉ lệ, $R_t$ là ma trận quay và $d_t$ là vector tịnh tiến. Sau khi căn chỉnh, sai khác giữa hai view tại frame $t$ và khớp $j$ có thể được tính bằng khoảng cách Euclid:

\begin{equation}
D_{t,j}
=
\left\|
\tilde{X}^{(1)}_{t,j}
-
\tilde{X}^{(2)}_{t,j}
\right\|_2.
\end{equation}

Giá trị $D_{t,j}$ càng lớn thì mức độ bất nhất giữa hai kết quả pose càng cao. Trong đồ án, đại lượng này được sử dụng như tín hiệu phát hiện sai lệch cục bộ giữa hai view, thay vì được xem là kết quả triangulation đầy đủ từ ảnh nhiều camera \cite{probtriangulation2023,gu2024multiview}.

Ràng buộc nhất quán đa góc nhìn có thể được biểu diễn dưới dạng hàm mất mát:

\begin{equation}
\mathcal{L}_{view}
=
\sum_{t=1}^{T}
\sum_{j=1}^{J}
w_{t,j}
\left\|
\tilde{X}^{(1)}_{t,j}
-
\tilde{X}^{(2)}_{t,j}
\right\|_2^2.
\end{equation}

Trong đó, $w_{t,j}$ là trọng số độ tin cậy ứng với từng khớp và từng frame. Trọng số này cần thiết vì độ tin cậy của từng khớp có thể thay đổi theo view, đặc biệt khi một khớp bị che khuất ở camera này nhưng được quan sát rõ hơn ở camera còn lại \cite{probtriangulation2023,gu2024multiview}.

\section{Phân tích cơ chế phân xử và hợp nhất pose}

Sau khi tính sai khác giữa hai view, hệ thống cần phân xử nguồn thông tin đáng tin cậy hơn cho từng khớp hoặc từng frame. Cách xử lý cục bộ theo khớp phù hợp vì lỗi monocular thường không xuất hiện đồng đều trên toàn bộ cơ thể, mà có thể tập trung tại các khớp bị che khuất hoặc có chuyển động khó quan sát \cite{smplify2016,physcap2020}. Nếu hai kết quả pose tương đối nhất quán, hệ thống có thể hợp nhất chúng bằng trung bình có trọng số. Nếu hai kết quả mâu thuẫn mạnh, hệ thống cần ưu tiên view có độ tin cậy cao hơn.

Gọi $C^{(1)}_{t,j}$ và $C^{(2)}_{t,j}$ lần lượt là điểm tin cậy của khớp $j$ tại frame $t$ trong view thứ nhất và view thứ hai. Công thức fusion theo trọng số có thể được viết như sau:

\begin{equation}
X^{F}_{t,j}
=
\frac{
C^{(1)}_{t,j}\tilde{X}^{(1)}_{t,j}
+
C^{(2)}_{t,j}\tilde{X}^{(2)}_{t,j}
}{
C^{(1)}_{t,j}
+
C^{(2)}_{t,j}
+
\varepsilon
}.
\end{equation}

Trong đó, $X^{F}_{t,j}$ là pose sau fusion, còn $\varepsilon$ là hằng số nhỏ để tránh chia cho $0$. Công thức này cho phép kết hợp hai nguồn thông tin theo độ tin cậy cục bộ, tương thích với các hướng multi-view sử dụng consistency hoặc weighted fusion để khai thác thông tin từ nhiều quan sát \cite{gu2024multiview,probtriangulation2023}.

Điểm tin cậy có thể được xây dựng từ nhiều tín hiệu khác nhau. Các tín hiệu này gồm sai khác đa góc nhìn, độ ổn định theo thời gian, biến thiên chiều dài xương và mức độ hợp lệ của pose theo ràng buộc động học \cite{physcap2020,delp2007opensim,uhlrich2023opencap}. Khi một khớp có sai khác đa góc nhìn lớn và đồng thời gây biến thiên bất thường trong chuỗi thời gian, hệ thống có cơ sở để xem khớp đó là ứng viên lỗi cần hiệu chỉnh mạnh hơn \cite{gu2024multiview,physcap2020}.

\section{Phân tích ràng buộc động học và cơ sinh học}

Pose 3D hợp lệ không chỉ cần phù hợp với dữ liệu quan sát mà còn cần phù hợp với cấu trúc cơ thể người. Các lỗi như chiều dài xương thay đổi bất thường, khớp xoay vượt giới hạn, foot skating hoặc floor penetration đã được chỉ ra là vấn đề thường gặp trong các hệ thống motion capture monocular nếu thiếu ràng buộc vật lý và cơ sinh học \cite{physcap2020}. Vì vậy, các ràng buộc động học và cơ sinh học được sử dụng để tăng tính hợp lý của pose sau refinement \cite{delp2007opensim,uhlrich2023opencap}.

Một ràng buộc cơ bản là bảo toàn chiều dài xương. Gọi $(a,b)$ là cạnh trong skeleton nối giữa khớp $a$ và khớp $b$. Chiều dài xương tại frame $t$ được tính như sau:

\begin{equation}
l_{t}^{(a,b)}
=
\left\|
X_{t,a} - X_{t,b}
\right\|_2.
\end{equation}

Nếu $l_{0}^{(a,b)}$ là chiều dài xương tham chiếu, ràng buộc bảo toàn chiều dài xương có thể được định nghĩa như sau:

\begin{equation}
\mathcal{L}_{bone}
=
\sum_{t=1}^{T}
\sum_{(a,b)\in\mathcal{B}}
\left(
l_{t}^{(a,b)}
-
l_{0}^{(a,b)}
\right)^2.
\end{equation}

Trong đó, $\mathcal{B}$ là tập các xương trong skeleton. Ràng buộc này giúp hạn chế biến dạng skeleton vì chiều dài các đoạn xương của cùng một người gần như không đổi trong một chuỗi chuyển động \cite{physcap2020,delp2007opensim}.

Ngoài chiều dài xương, giới hạn góc khớp cũng có thể được dùng để hạn chế các tư thế vượt miền chuyển động hợp lệ. Với một khớp có góc quay $\alpha$ và khoảng chuyển động hợp lệ $[\alpha_{min},\alpha_{max}]$, hàm phạt có thể được viết như sau:

\begin{equation}
\mathcal{L}_{angle}
=
\max(0, \alpha - \alpha_{max})^2
+
\max(0, \alpha_{min} - \alpha)^2.
\end{equation}

Hàm phạt này chỉ tạo giá trị khác $0$ khi góc khớp vượt ra ngoài khoảng hợp lệ. Cách xây dựng này giúp hệ thống không can thiệp mạnh vào các pose hợp lệ, đồng thời hạn chế các pose bất thường về mặt động học \cite{delp2007opensim,uhlrich2023opencap}. Tổng quát, thành phần ràng buộc cơ sinh học có thể được viết như sau:

\begin{equation}
\mathcal{L}_{bio}
=
\lambda_{bone}\mathcal{L}_{bone}
+
\lambda_{angle}\mathcal{L}_{angle}.
\end{equation}

Trong đồ án, ràng buộc cơ sinh học không nhằm thay thế một hệ mô phỏng cơ xương hoàn chỉnh. Vai trò chính của nó là cung cấp tiêu chí bổ sung để phát hiện và hạn chế các pose không hợp lý trong pipeline refinement \cite{delp2007opensim,uhlrich2023opencap,physcap2020}.

\section{Phân tích vai trò của SMPL và Learnable-SMPLify}

SMPL là mô hình cơ thể người tham số dựa trên skinning và blend shapes, cho phép biểu diễn hình dạng và tư thế cơ thể người trong một không gian có cấu trúc \cite{smpl2015}. Trong bài toán refinement, SMPL có vai trò quan trọng vì nó giúp ràng buộc các khớp và mesh vertices theo cấu trúc cơ thể người, thay vì xử lý các điểm 3D rời rạc độc lập \cite{smpl2015}. SMPLify sử dụng tối ưu lặp để fit tham số SMPL vào tín hiệu quan sát, từ đó tạo ra ước lượng pose và shape phù hợp hơn với mô hình cơ thể người \cite{smplify2016}.

Learnable-SMPLify thay thế quá trình fitting lặp của SMPLify bằng một mô hình neural regression nhằm giảm chi phí tối ưu và hỗ trợ vai trò plug-in post-processing cho các estimator sẵn có \cite{yang2025learnable}. Trong ngữ cảnh của đồ án, Learnable-SMPLify được xem là một thành phần hậu xử lý sau khi pose đã được đồng bộ, căn chỉnh và hợp nhất từ hai view. Cách sử dụng này phù hợp với định hướng của đồ án vì mục tiêu là refinement output có sẵn, không phải xây dựng mô hình sinh pose trực tiếp từ ảnh RGB \cite{yang2025learnable,shin2024wham}.

\section{Phân tích độ phức tạp tính toán của pipeline}

Độ phức tạp tính toán của pipeline phụ thuộc vào số lượng frame, số lượng khớp, số lượng vertices và số vòng lặp tối ưu. Gọi $T$ là số lượng frame, $J$ là số lượng khớp, $V$ là số lượng vertices của mesh SMPL và $K$ là số vòng lặp tối ưu hoặc số bước cập nhật trong refinement. Các bước xử lý trực tiếp trên keypoints 3D thường có chi phí thấp hơn so với các bước cần sinh lại mesh hoặc joints từ tham số SMPL \cite{smpl2015,smplify2016}.

Bước nạp dữ liệu và kiểm tra đầu vào có độ phức tạp xấp xỉ $O(TJ)$ nếu mỗi frame chứa $J$ khớp. Bước đồng bộ thời gian có độ phức tạp $O(T)$ khi hai chuỗi pose được đồng bộ theo chỉ số frame hoặc offset đơn giản. Bước căn chỉnh pose có độ phức tạp xấp xỉ $O(TJ)$ nếu similarity alignment được thực hiện trên tập khớp theo từng frame \cite{gu2024multiview}. Bước kiểm tra nhất quán đa góc nhìn và fusion đều có độ phức tạp $O(TJ)$ vì các phép tính được thực hiện theo từng khớp và từng frame \cite{gu2024multiview,probtriangulation2023}.

Nếu refinement được thực hiện trong không gian SMPL và mỗi vòng lặp cần sinh lại joints hoặc mesh vertices, độ phức tạp có thể xấp xỉ:

\begin{equation}
O(KT(J+V)).
\end{equation}

Thành phần $O(KT(J+V))$ thường chiếm chi phí lớn hơn trong pipeline nếu quá trình refinement trên SMPL được áp dụng cho toàn bộ chuỗi frame. Do đó, pipeline của đồ án ưu tiên thực hiện phát hiện sai lệch, phân xử và fusion trên keypoints 3D trước, sau đó mới sử dụng SMPL hoặc Learnable-SMPLify ở các bước cần bảo đảm tính hợp lệ của cơ thể người \cite{smpl2015,smplify2016,yang2025learnable}.

Bảng \ref{tab:complexity_analysis} tóm tắt độ phức tạp lý thuyết của các thành phần chính trong pipeline, với giả thiết chi phí sinh pose ban đầu từ WHAM được xem là đầu vào có sẵn và không tính vào pipeline hậu xử lý.

\begin{table}[H]
\centering
\caption{Phân tích độ phức tạp tính toán của các thành phần trong pipeline}
\label{tab:complexity_analysis}
\begin{tabular}{|p{4.2cm}|p{4.2cm}|p{4.8cm}|}
\hline
\textbf{Thành phần} & \textbf{Độ phức tạp xấp xỉ} & \textbf{Giả thiết} \\
\hline
Đọc và kiểm tra joints 3D & $O(TJ)$ & Chỉ xét tọa độ khớp, không đọc toàn bộ mesh. \\
\hline
Đọc vertices/mesh nếu sử dụng & $O(TV)$ & Chỉ phát sinh khi pipeline dùng SMPL mesh. \\
\hline
Đồng bộ theo chỉ số frame/timestamp & $O(T)$ & Áp dụng khi có timestamp hoặc hai video cùng FPS. \\
\hline
Đồng bộ bằng tín hiệu chuyển động & $O(WT)$ hoặc $O(T \log T)$ & Phụ thuộc cách tìm offset/correlation. \\
\hline
Căn chỉnh pose & $O(TJ)$ & Ước lượng similarity/Procrustes trên $J$ khớp. \\
\hline
Kiểm tra nhất quán & $O(TJ)$ & Tính sai khác theo từng frame và khớp. \\
\hline
Phân xử và fusion & $O(TJ)$ & Tính trọng số tin cậy và hợp nhất pose. \\
\hline
Refinement trên keypoints & $O(KTJ)$ & $K$ là số vòng lặp hoặc số lần cập nhật. \\
\hline
Refinement trong không gian SMPL & $O(KT(J+V))$ & Phụ thuộc số vòng lặp, số khớp và số đỉnh mesh. \\
\hline
\end{tabular}
\end{table}

Nếu bỏ qua chi phí sinh pose ban đầu từ WHAM, vì WHAM được xem là mô hình đầu vào có sẵn, độ phức tạp chính của pipeline hậu xử lý có thể mô tả ở mức:

\begin{equation}
O(TJ) + O(KT(J+V)).
\end{equation}

Biểu thức trên cho thấy các bước xử lý keypoints có chi phí tuyến tính theo số frame và số khớp, trong khi bước refinement trong không gian SMPL phụ thuộc thêm vào số vertices và số vòng lặp tối ưu \cite{smpl2015,smplify2016,yang2025learnable}.

\section{Nhận xét về vai trò của các thành phần trong pipeline}

Mỗi thành phần trong pipeline xử lý một nhóm lỗi khác nhau của output monocular. Ràng buộc dữ liệu giúp hạn chế hiệu chỉnh quá mức so với output ban đầu \cite{smplify2016,shin2024wham}. Ràng buộc thời gian giúp giảm jitter và các biến đổi bất thường giữa những frame liên tiếp \cite{physcap2020}. Ràng buộc đa góc nhìn hỗ trợ phát hiện sai lệch giữa hai kết quả monocular độc lập \cite{gu2024multiview,probtriangulation2023}. Ràng buộc động học và cơ sinh học giúp hạn chế các pose không hợp lý về cấu trúc cơ thể \cite{delp2007opensim,uhlrich2023opencap,physcap2020}. Ràng buộc SMPL giúp đưa kết quả về không gian biểu diễn cơ thể người có cấu trúc \cite{smpl2015,smplify2016}.

Bảng \ref{tab:constraint_roles} tóm tắt vai trò của các thành phần ràng buộc trong phương pháp đề xuất và nêu thêm rủi ro nếu cấu hình không phù hợp.

\begin{table}[H]
\centering
\caption{Vai trò của các thành phần ràng buộc trong phương pháp đề xuất}
\label{tab:constraint_roles}
\begin{tabular}{|p{3.4cm}|p{4.1cm}|p{3.0cm}|p{3.6cm}|}
\hline
\textbf{Thành phần ràng buộc} & \textbf{Vai trò} & \textbf{Loại lỗi xử lý} & \textbf{Rủi ro nếu cấu hình sai} \\
\hline
Ràng buộc dữ liệu & Giữ pose sau hiệu chỉnh gần với output ban đầu. & Tránh hiệu chỉnh quá mức. & Nếu quá mạnh, hệ thống khó sửa lỗi thật. \\
\hline
Ràng buộc thời gian & Làm mượt chuyển động. & Jitter, biến đổi đột ngột. & Nếu quá mạnh, làm mất chuyển động nhanh. \\
\hline
Ràng buộc đa góc nhìn & Phát hiện mâu thuẫn giữa hai view. & Sai lệch chiều sâu, lỗi cục bộ theo view. & Nếu đồng bộ/căn chỉnh sai, có thể đánh dấu nhầm. \\
\hline
Ràng buộc chiều dài xương & Duy trì skeleton ổn định. & Biến dạng xương, thay đổi chiều dài bất thường. & Nếu quá cứng, giảm tính tự nhiên của chuyển động. \\
\hline
Ràng buộc góc khớp & Hạn chế tư thế vượt miền hợp lệ. & Pose không hợp lý về động học. & Có thể loại nhầm tư thế đặc biệt nhưng hợp lệ. \\
\hline
Ràng buộc SMPL & Đưa pose về biểu diễn cơ thể người có cấu trúc. & Thiếu nhất quán giữa joints và mesh. & Tăng chi phí tính toán, phụ thuộc chất lượng khởi tạo. \\
\hline
\end{tabular}
\end{table}

Sự kết hợp các thành phần trên phù hợp với đặc điểm của bài toán hậu xử lý vì mỗi thành phần xử lý một loại sai số khác nhau. Tuy nhiên, các trọng số cần được lựa chọn dựa trên thực nghiệm, vì ràng buộc thời gian quá mạnh có thể làm mất chuyển động nhanh, ràng buộc đa góc nhìn có thể gây lỗi nếu đồng bộ chưa tốt, và ràng buộc cơ sinh học quá cứng có thể làm giảm sự đa dạng của các tư thế hợp lệ \cite{physcap2020,gu2024multiview,delp2007opensim}.

\section{Kết chương}

Chương này đã trình bày phân tích lý thuyết cho phương pháp đề xuất trong đồ án. Bài toán refinement được mô hình hóa dưới dạng tối ưu nhiều thành phần, trong đó các ràng buộc dữ liệu, thời gian, đa góc nhìn, động học/cơ sinh học và SMPL cùng tham gia vào quá trình hiệu chỉnh pose \cite{smpl2015,smplify2016,physcap2020,gu2024multiview}. Các phân tích cho thấy phương pháp đề xuất phù hợp với bài toán hậu xử lý kết quả monocular pose estimation, vì hệ thống không thay thế mô hình ước lượng ban đầu mà bổ sung cơ chế kiểm tra và hiệu chỉnh sau ước lượng \cite{shin2024wham}.

Chương này cũng đã phân tích độ phức tạp tính toán của các giai đoạn chính trong pipeline. Kết quả phân tích cho thấy các bước xử lý trên keypoints 3D có chi phí tuyến tính theo số frame và số khớp, trong khi các bước refinement trong không gian SMPL có chi phí cao hơn do phụ thuộc vào số vertices và số vòng lặp tối ưu \cite{smpl2015,smplify2016,yang2025learnable}. Đây là cơ sở để thiết kế thực nghiệm đánh giá hiệu quả và chi phí xử lý của phương pháp trong chương tiếp theo.

\end{document}
